\documentclass[DIN, pagenumber=false, fontsize=11pt, parskip=half]{scrartcl}

\usepackage[utf8]{inputenc}
\usepackage[T1]{fontenc}
\usepackage{template}

\titlehead{\centering\small
    Worked solutions to selected exercises from\\
    Richard Hammack, \emph{Book of Proof, edition 2.2}\\
\rule{0.4\linewidth}{0.4pt}}

\title{}
\author{} % Oscar A. Carrera
\date{}

\begin{document}
\maketitle

\setcounter{section}{4}

\section{Contrapositive Proof}

\textbf{A. Use the method of contrapositive proof to prove the following statements.}

\begin{problems}
    \problem Suppose $n \in \mathbb{Z}$. If $n^2$ is odd, then $n$ is odd.
    \begin{proof}
        We prove the contrapositive: if $n$ is not odd, then $n^2$ is not odd.\\
        Suppose that $n$ is not odd.
        Thus $n$ is even, so $n = 2a$ for some integer $a$.
        Consequently $n^2 = (2a)^2 = 4a^2 = 2(2a^2).$
        Therefore $n^2 = 2b$, where the integer $b = 2a^2$.
        Consequently, $n^2$ is even by definition of an even number.
        Therefore $n^2$ is not odd.
    \end{proof}

    \problem Suppose $a,b,c \in \mathbb{Z}$. If $a$ does not divide $bc$, then $a$ does not divide $b$.
    \begin{proof}
        We prove the contrapositive: if $a \mid b$, then $a \mid bc$.\\
        Suppose $a \mid b$.
        Then $b = ak$ for some $k \in \mathbb{Z}$.
        Multiplying both sides by $c$ gives $bc = a(kc)$.
        Since $kc \in \mathbb{Z}$, it follows that $a \mid bc$.
    \end{proof}

    \problem Suppose $x \in \mathbb{R}$. If $x^3-x > 0$ then $x > -1$.
    \begin{proof}
        We prove the contrapositive: if $x \leq -1$ then $x^3-x \leq 0$.\\
        Suppose $x \leq -1$. Then $x < 0$, $x - 1 < 0$, and $x + 1 \leq 0$.
        Notice that $x^3 - x = x(x-1)(x+1)$, a product of two negative factors and one nonpositive factor.
        Since the product of two negative numbers and a nonpositive number is nonpositive, we have $x(x-1)(x+1) \leq 0$, that is, $x^3 - x \leq 0$.
    \end{proof}
\end{problems}

\end{document}
